Nennen und erläutern Sie wesentliche Kriterien zur Abgrenzung von
Privatversicherung und Sozialversicherung.

INSERT FIGURE PAGE 5 DAV

Träger der GKV sind die Krankenkassen als selbstverwaltete öffentlich rechtli-
che Körperschaften, GKV unterliegen insoweit weder den allgemeinen privatrechtlichen Vorschriften
des BGB's, HGB's, VAG oder VVG's. Sie haben keine handelsrechtliche
Buchführungspflicht und ebenso keine Körperschaftssteuerpflicht.

Träger der PKV sind private Krankenversicherungsunternehmen, die den all-
gemeinen und speziellen Normen des Zivilrechts sowie dem Versicherungsauf-
sichtsgesetz (VAG) unterliegen.


Erläutern Sie das grundlegende Prinzip der Prämienkalkulation in der
privaten Krankenversicherung.

Zum Vertragsabschlusszeitpunkt schätzt der Versicherer für jeden einzelnen Ver-
sicherten, den Barwert der lebenslang erwarteten Krankheitskosten. Dem Prinzip
der individuellen Risikoäquivalenz folgend, muss für jeden Vertrag zum Vertragsabschlusszeitpunkt 
der Barwert der erwarteten Krankheitskosten dem Barwert aller erwarteten Prämienzahlungen entsprechen.
Rechnungsgrundlagen sind in der Krankenversicherungsauf-
sichtsverordnung (KVAV) festgelegt und umfassen im Wesentlichen die sog. Kopf-
schäden, den Rechnungszins, die Ausscheideordnung (Sterblichkeit und Storno),
den Sicherheitszuschlag, weitere Zuschläge zur Deckung von Abschluss-, Verwal-
tungs- und Regulierungskosten sowie Zuschläge für die erfolgsunabhängige RfB.

Die jährlich einzelvertraglich erwarteten Krankheitskosten (Kopfschäden) steigen mit zunehmen-
dem Lebensalter. 
Eine an dem natürlichen Krankheitskostenverlauf angepasste Prämienzahlung würde dement-
sprechend ebenfalls periodisch ansteigen (natürliche Prämie). 
Dem Kalkulationsprinzip der Lebensversicherung folgend wird jedoch eine periodisch konstante
Prämienzahlung vereinbart (tatsächliche Prämie)
In der Regel sind daher die zu Vertragsbeginn vom Versicherten
erwartungsgemäß geleisteten jährlichen Prämienzahlungen größer als die jährlich
erwarteten Leistungszahlungen. (Altersrückstellungen)



Aus welchen Bestandteilen setzt sich das 3-Schichtenmodell der Alters-
vorsorge in Deutschland zusammen?

1. Schicht: Gesetzliche Vorsorge
Gesetzliche Rentenversicherung, Beamtenversorgung, Berufsständische
Versorgung, Alterssicherung für Landwirte, Künstlersozialversicherung
2. Schicht: Erwerbsbasierte Vorsorge
Staatlich geförderte, privatrechtliche Verträge, insb. betriebliche Altersvorsorge,
„Rürup“-Rente (häufig der ersten Schicht zugerechnet) und „Riester“-Rente
3. Schicht: Private Vorsorge
Bildung von Finanzvermögen, z.B. Wertpapiere, Sparguthaben,
Kapital- und Rentenversicherung, Bildung von Sachvermögen, z.B. Immobilien


Vom Beitragszahler der GRV (1. Schicht) werden die Beiträge nicht angespart, son-
dern im Rahmen eines Umlageverfahrens als Rentenleistungen an die gegen-
wärtigen Rentner unmittelbar ausbezahlt („pay as you go“-System). Beitragszahler 
haben demnach auch keinen Rückzahlungsanspruch auf die geleisteten Beiträge, 
sehr wohl aber eine Anwartschaft auf Rentenleistungen.
Es ergibt sich der sog. Generationenvertrag bzw. das Solidaritätsprinzip.

In der betrieblichen Altersvorsorge (bAV) (2. Schicht) erteilt der Arbeitgeber dem Arbeit-
nehmer aus Anlass des Arbeitsverhältnisses eine Versorgungszusage.
Heute haben Arbeitnehmer grundsätzlich einen Anspruch auf eine bAV (Betriebsrentengesetz),
durch Entgeltumwandlung, also Verzicht auf einen Teil ihres Lohnes bzw. Gehalts. 
Steuerliche Anreize und staatliche Zulagen-Förderung steigern die Attraktivität.
Sie kann auf fünf verschiedene Arten durchgeführt werden: 
- Als Direktzusage durch den Arbeitgeber 
- Direktversicherung durch den Arbeitgeber 
- Pensionskasse
- Pensionsfonds 
- Unterstützungskasse.
Zusätzlich können Lebensarbeitszeitkonten, auf denen ein Arbeitnehmer Überstunden einzahlen, genutzt werden.

Probleme sind die Portabilität (z.B. Wechsel von Arbeitgeber), der Unverfallbarkeit von Ansprüchen der Arbeitnehmer, 
aber auch die Durchgriffshaftung bei externen Durchführungswegen.

Riester-Rente \& Rürup-Rente (2. Schicht)
Die Riester-Rente ist eine freiwillige Versicherung. 
Ihre Förderung zielt auf rentenversicherungspflichtige Arbeitnehmer und
Selbständige und ist durch das Altersvermögensgesetz geregelt
Rürup-Rente ist ebenfalls eine freiwillige Versicherung. Sie richtet
sich an nicht rentenversicherungspflichtige Selbständige, Freiberufler und Beamte.
Sie unterliegen als geförderte Vorsorgemodelle besonderer Normen, wie etwa dem 
Beleihungs-, Verpfändungs- und Veräußerungsverbot.



Beschreiben Sie warum man von der „doppelten Alterung der Gesell-
schaft“ und erläutern Sie den Einfluss des demografischen Wandels auf
die gesetzliche Rentenversicherung (GRV) hat.

Als Demografischen Wandel werden Tendenzen der Bevölkerungsentwicklung
bezeichnet, wie z.B. Alters/Migranten/Geschlecht-struktur der Bevölkerung und die Entwicklung der Gebur-
ten- und Sterbezahlen umfassen.
Es wird in Deutschland das Phänomen der doppelten Alterung der Gesellschaft beschreiben, 
d.h. es werden weniger Menschen geboren und zugleich leben die Menschen immer länger.
Wer im Jahr 1964 geboren ist, wird nach 67 Jahren die gesetzliche Altersrente empfangen.

Die Vorteile der GRV sind die kapitalgedeckten Finanzierungssysteme, welche auf 
auf einen stabilen Geldwert und eine angemessene Kapitalmarktverzinsung angewiesen sind.
In einer stagnierenden oder gar schrumpfenden Wirtschaft wird die Erwirtschaftung 
einer angemessenen Kapitalmarktverzinsung erschwert. Kapitalbildende Vorsorgemodelle 
sind dann auf offene Volkswirtschaften mit stabilen Währungsbeziehungen angewiesen.
Nun ist aber die Bevölkerungsentwicklung grundsätzlich ein wesentlicher Treiber
des Wirtschaftswachstums. Der demografische Wandel hat daher sowohl quantitativen 
wie auch strukturellen Einfluss auf das Wirtschaftswachstum und das
Kaufverhalten.



